% !TEX root = chapter-standalone.tex

\section{Opposite of a preorder}\label{sec:opposite-of-a-poset}

\todotext{Add material explaining / commenting about dualities, e.g. between meet and join}

\begin{ctdefinition}[Opposite of a preorder]
    \label{def:poset-opposite}
    The \maindef{opposite of a preorder}~$\posAdefinition$ is the \SY{preorder} denoted~$\posAop=\tupp{\posAset, \posAleqop}$.
    It has the same elements as~\posA, but is equipped with the reverse ordering, in the sense that, for all~$\posela,\poselb \setin \posAset$,
    \begin{equation}\label{eq:poset-opposite}
        \prfdoubleperiod{
            \posela \posAleq \poselb
        }{
            \poselb \posAleqop \posela
        }
    \end{equation}
\end{ctdefinition}

\begin{lemma}
    \label{lem:opposite-preorder}
    Let~\posA be a \SY{preorder}.
    Then the relation~$\posAleqop$ is reflexive and transitive, so~$\posAop$ is a \SY{preorder}.
    If, moreover,~\posA is a \SY{poset}, then~$\posAop$ is a \SY{poset}.
\end{lemma}
\begin{proof}
    Each property of~$\posAleqop$ follows from the same property of~$\posAleq$, read backwards.
    \begin{enumerate}
        \item Reflexivity: for any~$\posela \setin \posAset$, we have~$\posela \posAleq \posela$, and therefore~$\posela \posAleqop \posela$.
        \item Transitivity: suppose~$\posela \posAleqop \poselb$ and~$\poselb \posAleqop \poselc$.
        By definition, this means~$\poselb \posAleq \posela$ and~$\poselc \posAleq \poselb$.
        By transitivity of~$\posAleq$, we get~$\poselc \posAleq \posela$, that is,~$\posela \posAleqop \poselc$.
        \item Antisymmetry (when~\posA is a \SY{poset}): suppose~$\posela \posAleqop \poselb$ and~$\poselb \posAleqop \posela$.
        Then~$\poselb \posAleq \posela$ and~$\posela \posAleq \poselb$, so~$\posela = \poselb$ by antisymmetry of~$\posAleq$.
    \end{enumerate}
\end{proof}
For a given~$\posela \setin \posAset$, we will sometimes write~$\posela^*$ to denote its corresponding copy in~$\posAop$, in order to emphasize that~$\posela$ and~$\posela^*$ belong to distinct \SY{preorders}.
However, often we will not be so pedantic with our notation.

\begin{marginfigure}
    \centering
    \includesag{40_dpcatfig_opposite_a}\\
    \includesag{40_dpcatfig_opposite_b}
    \caption{
        Opposite of a preorder.
    }
    \label{fig:poset-opposite}
\end{marginfigure}

\begin{example}[Credit and debt]
    Let us define the set
    %
    \begin{equation}\label{eq:credit-debit}
        \posAset=\makeset{0.00,0.01,0.02,\dots}\setsubseteq\reals
    \end{equation}
    %
    of all \CHFneutral \ monetary quantities approximated to the cent.
    From this set we can define two \SY{posets},~$\textcolor{functionality}{\posgenA}^{+} = \tup{\posAset, {{\Rleq}}}$ and~$\textcolor{requirements}{\posgenA}^{-} = \tup{\posAset,{{ \Rgeq}}}$, that are the opposite of each other.
    If the context is that, given two quantities \unit[1]{\CHFneutral} and \unit[2]{\CHFneutral}, we prefer \unit[1]{\CHFneutral} to \unit[2]{\CHFneutral} (for example because it is a cost to pay to acquire a component), then we are working in~$\textcolor{functionality}{\posgenA}^{+}$, otherwise we are working in~$\textcolor{requirements}{\posgenA}^{-}$ (for example because it represents the price at which we are selling our product).
    Traditionally, in double-entry ledger systems, the numbers were not written with negative signs, but rather in color: red and black.
    From this convention we get the idioms ``being in the black'' and ``being in the red''.
\end{example}
